\documentclass[../main.tex]{subfiles}
\begin{document}

\section{Overview}

Chapter 2 established that attribution graphs expose the causal structure of a
single model prediction, but also showed that their size makes direct human
inspection difficult. Chapter 3 presents the proposed summarization method for
turning such an attribution graph into a smaller \emph{summary graph} whose nodes
represent groups of related features and whose edges preserve the main causal
flow. The method has three stages. \textbf{Pruning} removes attribution-graph
nodes and edges that carry little input-to-output influence. \textbf{Clustering}
groups the remaining feature nodes by signed role similarity, while keeping
embedding and logit nodes fixed as singletons. \textbf{Interpreting} aggregates
the clustered graph into a directed acyclic summary graph and assigns display
labels to its supernodes. Section~\ref{sec:problem} formalizes the input, output,
desiderata, and constraints of the summarization problem. Section~\ref{sec:method}
then defines the three stages of the method in sequence.
Figure~\ref{fig:methodology-pipeline} gives an overview of the pipeline on the
running ``Dallas $\to$ Austin'' example.

% Methodology pipeline figure.
% Requires (loaded in main.tex): tikz + libraries calc, positioning, arrows.meta, shapes.geometric.
\begin{figure}[t]
\centering
\resizebox{\textwidth}{!}{%
\begin{tikzpicture}[
  >=Stealth,
  emb/.style   ={draw, diamond, aspect=1.2, inner sep=0pt, minimum size=4.7mm,
                 fill=violet!23, draw=violet!65!black},
  feat/.style  ={circle, draw=black!85, line width=0.45pt, inner sep=0pt,
                 minimum size=3.7mm,
                 fill=white},
  faint/.style ={circle, draw=black!22, inner sep=0pt, minimum size=3.6mm,
                 fill=black!4},
  err/.style   ={regular polygon, regular polygon sides=3, draw=red!55!black,
                 line width=0.45pt, inner sep=0pt, minimum size=4.4mm,
                 fill=red!18},
  logit/.style ={draw, rounded corners=1pt, inner sep=1.8pt, minimum width=9mm,
                 minimum height=5.3mm, fill=black!82, text=white, font=\tiny},
  sn/.style    ={ellipse, draw=teal!70!black, line width=0.75pt, inner sep=1.7pt,
                 minimum width=14mm, minimum height=8.7mm, fill=cyan!15,
                 font=\scriptsize},
  labelnode/.style={ellipse, draw=teal!70!black, line width=0.75pt, inner sep=2pt,
                 minimum width=15mm, minimum height=9.5mm, fill=cyan!15,
                 font=\scriptsize, align=center},
  dot/.style   ={circle, fill=black!85, inner sep=0pt, minimum size=1.55mm},
  pale/.style  ={circle, fill=black!28, inner sep=0pt, minimum size=1.15mm},
  panel/.style ={draw=black!52, line width=0.5pt, rounded corners=3pt, fill=gray!3,
                 minimum width=33mm, minimum height=39mm},
  parrow/.style={->, line width=1.25pt, gray!70!black},
  kept/.style  ={->, black!88, line width=0.75pt},
  denseedge/.style={->, gray!58, line width=0.36pt},
  dropped/.style={->, gray!32, line width=0.28pt, dashed},
  cycle/.style ={->, orange!90!black, line width=1.05pt},
  removed/.style={orange!90!black, line width=0.85pt, dashed},
  ecap/.style  ={font=\scriptsize, align=center, text width=31mm},
  slab/.style  ={font=\small\bfseries, align=center},
  note/.style  ={font=\tiny, align=center, text=orange!65!black},
]

% ---------- Panel A: dense attribution graph ----------
\begin{scope}[shift={(0,0)}]
  \node[panel] (pA) {};
  \node[emb]   (aEmb1) at (-0.95,-1.45) {};
  \node[emb]   (aEmb2) at ( 0.00,-1.45) {};
  \node[emb]   (aEmb3) at ( 0.95,-1.45) {};
  \node[err]   (aErr1) at (-1.28,-0.08) {};
  \node[err]   (aErr2) at ( 1.20, 0.70) {};
  \node[feat]  (a1) at (-0.98,-0.55) {};
  \node[feat]  (a2) at (-0.35,-0.55) {};
  \node[feat]  (a3) at ( 0.30,-0.55) {};
  \node[feat]  (a4) at ( 0.95,-0.55) {};
  \node[feat]  (a5) at (-0.78, 0.28) {};
  \node[feat]  (a6) at (-0.13, 0.28) {};
  \node[feat]  (a7) at ( 0.52, 0.28) {};
  \node[feat]  (a8) at ( 0.05, 0.92) {};
  \node[logit] (aLogit) at ( 0.00, 1.52) {Austin};

  \foreach \s/\t in {aEmb1/a1,aEmb1/a2,aEmb1/a5,aEmb2/a1,aEmb2/a2,
                     aEmb2/a3,aEmb2/a6,aEmb3/a3,aEmb3/a4,aEmb3/a7,
                     aErr1/a1,aErr1/a5,aErr2/a7,aErr2/aLogit,
                     a1/a5,a1/a6,a2/a5,a2/a6,a2/a7,a3/a6,a3/a7,
                     a4/a7,a5/a8,a6/a8,a7/a8,a5/aLogit,a6/aLogit,
                     a7/aLogit,a8/aLogit}
     \draw[denseedge] (\s)--(\t);
\end{scope}

% ---------- Panel B: pruned graph ----------
\begin{scope}[shift={(4.85,0)}]
  \node[panel] (pB) {};
  \node[emb]   (bEmb1) at (-0.78,-1.45) {};
  \node[emb]   (bEmb2) at ( 0.58,-1.45) {};
  \node[faint] (bGone1) at ( 0.00,-0.95) {};
  \node[faint] (bGone2) at ( 1.02, 0.36) {};
  \node[err, draw=red!20, fill=red!6] (bErr) at (-1.13,-0.25) {};
  \node[feat]  (b1) at (-0.60,-0.48) {};
  \node[feat]  (b2) at ( 0.42,-0.48) {};
  \node[feat]  (b3) at (-0.05, 0.45) {};
  \node[logit] (bLogit) at ( 0.00, 1.52) {Austin};

  \foreach \s/\t in {bEmb1/b1,bEmb2/b2,b1/b3,b2/b3,b3/bLogit,b2/bLogit}
     \draw[kept] (\s)--(\t);
  \foreach \s/\t in {bErr/b1,bGone1/b3,bGone2/bLogit}
     \draw[dropped] (\s)--(\t);
\end{scope}

% ---------- Panel C: clustered graph before cycle resolution ----------
\begin{scope}[shift={(9.7,0)}]
  \node[panel] (pC) {};
  \node[emb]   (cEmb1) at (-0.82,-1.45) {};
  \node[emb]   (cEmb2) at ( 0.62,-1.45) {};
  \node[sn]    (cSN1) at (-0.66,-0.38) {};
  \node[dot]   at (-0.88,-0.38) {}; \node[dot] at (-0.63,-0.20) {};
  \node[dot]   at (-0.43,-0.43) {};
  \node[sn]    (cSN2) at ( 0.66,-0.05) {};
  \node[dot]   at ( 0.44,-0.05) {}; \node[dot] at ( 0.69, 0.13) {};
  \node[dot]   at ( 0.89,-0.10) {};
  \node[sn]    (cSN3) at ( 0.02, 0.92) {};
  \node[dot]   at (-0.17, 0.92) {}; \node[dot] at ( 0.20, 0.92) {};
  \node[logit] (cLogit) at ( 0.00, 1.54) {Austin};

  \draw[kept] (cEmb1)--(cSN1);
  \draw[kept] (cEmb2)--(cSN2);
  \draw[cycle] (cSN1) to[bend left=15] (cSN2);
  \draw[cycle] (cSN2) to[bend left=15] (cSN1);
  \draw[kept] (cSN1)--(cSN3);
  \draw[kept] (cSN2)--(cSN3);
  \draw[kept] (cSN3)--(cLogit);
  \node[note] at (0.05, -0.85) {cycle found};
\end{scope}

% ---------- Panel D: resolved and labeled summary ----------
\begin{scope}[shift={(14.55,0)}]
  \node[panel] (pD) {};
  \node[emb,label={[font=\tiny,text=violet!65!black]below:Dallas}]  (dEmb1) at (-0.82,-1.42) {};
  \node[emb,label={[font=\tiny,text=violet!65!black]below:capital}] (dEmb2) at ( 0.82,-1.42) {};
  \node[labelnode] (dState) at (-0.78,-0.35) {Texas\\state};
  \node[labelnode] (dCapital) at (0.78,-0.35) {capital\\query};
  \node[labelnode] (dAnswer) at (0.02,0.78) {answer\\selection};
  \node[logit] (dLogit) at (0.00,1.55) {Austin};

  \draw[kept] (dEmb1)--(dState);
  \draw[kept] (dEmb2)--(dCapital);
  \draw[kept] (dState)--(dAnswer);
  \draw[kept] (dCapital)--(dAnswer);
  \draw[kept] (dAnswer)--(dLogit);
  \draw[removed] (dCapital) to[bend left=20] (dState);
  \node[note] at (0.08, -0.87) {cycle edge removed};
\end{scope}

% ---------- pipeline arrows ----------
\draw[parrow] (pA.east)--(pB.west);
\draw[parrow] (pB.east)--(pC.west);
\draw[parrow] (pC.east)--(pD.west);

% ---------- column headers ----------
\node[slab, above=5mm of pA.north] {Dense graph};
\node[slab, above=5mm of pB.north] {Pruned graph};
\node[slab, above=5mm of pC.north] {Clustered graph};
\node[slab, above=5mm of pD.north] {Final summary};

% ---------- per-panel captions ----------
\node[ecap, below=5.5mm of pA.south] {Many feature and error nodes are connected by dense causal pathways.};
\node[ecap, below=5.5mm of pB.south] {Weak pathways and error nodes are removed, leaving the main route to the output.};
\node[ecap, below=5.5mm of pC.south] {Related feature nodes are grouped, but aggregated links may introduce cycles.};
\node[ecap, below=5.5mm of pD.south] {Cycle-causing links are removed and each supernode receives a readable label.};

% ---------- legend ----------
\begin{scope}[shift={(0,-5.35)}]
  \node[emb,scale=0.78] (g1) at (-0.65,0) {}; \node[right=1.5mm of g1, font=\scriptsize] {input token};
  \node[feat]           (g2) at ( 2.70,0) {}; \node[right=1.5mm of g2, font=\scriptsize] {feature};
  \node[err,scale=0.85] (g3) at ( 4.95,0) {}; \node[right=1.5mm of g3, font=\scriptsize] {error node};
  \node[logit]          (g4) at ( 7.75,0) {Austin}; \node[right=1.5mm of g4, font=\scriptsize] {output};
  \node[sn,scale=0.76]  (g5) at (10.55,0) {}; \node[right=1.5mm of g5, font=\scriptsize] {supernode};
  \draw[removed] (13.35,-0.05)--(14.05,-0.05); \node[font=\scriptsize] at (14.95,-0.05) {removed cycle edge};
\end{scope}

\end{tikzpicture}}
\caption[Methodology pipeline.]{Methodology pipeline. The method starts from a
dense attribution graph, prunes it to the main causal pathways, clusters related
feature nodes into supernodes, then removes cycle-causing links and assigns
readable labels so the final summary can be inspected as a causal story.}
\label{fig:methodology-pipeline}
\end{figure}


\section{Problem Formulation}
\label{sec:problem}

\subsection{Input}

The input is an attribution graph describing the steps a model used to produce an
output for a target token on a particular prompt. We write this graph as
$G = (V, E, W)$, where $V$ is the node set, $E$ is the directed edge set, and
$W \in \mathbb{R}^{N \times N}$ with $N := |V|$ is the signed weighted adjacency
matrix. The graph contains four types of nodes:
\[
V = V_{\text{emb}} \cup V_{\text{feature}} \cup V_{\text{error}} \cup
V_{\text{logit}}.
\]
Embedding nodes $V_{\text{emb}}$ represent input-token sources, feature nodes
$V_{\text{feature}}$ represent transcoder features, error nodes
$V_{\text{error}}$ represent reconstruction residuals, and logit nodes
$V_{\text{logit}}$ represent output predictions.

The edges in the graph represent linear effects between nodes. Under the matrix
convention used throughout this thesis, $W_{ij}$ is the signed effect from source
node $j$ to target node $i$, so columns are sources and rows are targets. The
magnitude $|W_{ij}|$ represents the strength of the influence, and the sign of
$W_{ij}$ represents whether the influence is positive or negative. For each
feature node, its activity is the sum of the signed contributions from its
incoming edges. Because these edges follow the forward computation of the model,
$G$ is a directed acyclic graph.

The input also contains two external weighting vectors that identify which parts
of the attribution graph should be emphasized. The token weight vector
$\tau \in \mathbb{R}^N$ identifies the input tokens of interest by assigning
nonzero weight to their embedding nodes. The logit weight vector
$\ell \in \mathbb{R}^N$ identifies the output predictions of interest by
assigning nonzero weight to their logit nodes. These vectors are not additional
model-computation nodes or edges; they are scoring inputs used by the pruning
stage.

\subsection{Output}

The output is a \emph{summary graph} $H = (V_H, E_H, W_H)$. Its node set
$V_H = \{C_1, C_2, \dots, C_K\}$ consists of \emph{supernodes}, where each
supernode is a nonempty subset of the non-error attribution-graph nodes:
\begin{equation}
V_H \subseteq \mathcal{P}\!\left(V \setminus V_{\text{error}}\right)
\setminus \{\emptyset\}.
\label{eq:summary-node-set}
\end{equation}
Writing $U_H := \bigcup_{C \in V_H} C$ for the attribution-graph nodes covered
by the summary, admissible supernodes satisfy
\begin{equation}
C_i \cap C_j = \emptyset \quad (i \ne j),
\qquad
\forall v \in U_H \cap (V_{\text{emb}} \cup V_{\text{logit}}),\; \{v\} \in V_H.
\label{eq:supernode-conditions}
\end{equation}
Therefore $V_H$ forms a partition of the represented subset
$U_H \subseteq V \setminus V_{\text{error}}$. Embedding and logit nodes are
represented only as singleton supernodes, so only feature nodes may be grouped.
The edge set $E_H$ and signed weights $W_H$ describe the net causal influence
between supernodes.

\subsection{Desiderata of the summary graph}

A useful summary graph must be small enough to inspect while still representing
the causal computation of the original attribution graph. Four desiderata make
this requirement explicit.
\begin{description}[leftmargin=2.0em,style=nextline,topsep=2pt,itemsep=3pt]
\item[(D1) Importance retention.] The summary should retain the nodes and edges
that lie on high-mass pathways between the seed embedding nodes and seed logit
nodes, and discard nodes and edges that do not contribute to those pathways.
\item[(D2) Atomicity.] Each supernode should represent one causal role. Its
member features should have similar signed downstream targets and similar signed
upstream sources, so the reader can interpret the supernode as a single
mechanism.
\item[(D3) Causal-information preservation.] Causal relations between retained
nodes should remain visible as edges between supernodes when possible. If two
directly connected nodes are merged into one supernode, their edge becomes
internal and no longer appears in $H$.
\item[(D4) Interpretability.] The final graph should be a directed acyclic graph
with few enough supernodes for manual inspection. Acyclicity is required because
the summary is intended to show a forward causal story.
\end{description}

\subsection{Constraints}
\label{sec:constraints}

Two hard constraints implement the structural part of interpretability. First,
$H$ must be acyclic:
\begin{equation}
H \text{ is a directed acyclic graph.}
\label{eq:acyclicity}
\end{equation}
Second, the number of supernodes must not exceed a complexity budget:
\begin{equation}
K = |V_H| \le K_{\max}.
\label{eq:budget}
\end{equation}
The method does not treat these constraints as soft losses. The clustering stage
enforces the budget, and the interpreting stage enforces acyclicity.

\section{Method}
\label{sec:method}

The method builds the summary graph in three sequential stages. Pruning first
chooses a retained attribution subgraph $(V', E')$ by applying retained-mass
thresholds to node and edge importance scores. Clustering then maps the retained
nodes to supernodes through a partition $f : V' \to V_H$ that optimizes atomicity
and causal-information preservation under the budget $K_{\max}$. Interpreting
finally converts that partition into a directed acyclic summary graph by
aggregating edge weights and removing cycles introduced by aggregation.

The method uses computable proxies because faithfulness and interpretability
cannot be measured directly while the graph is being constructed. Faithfulness is
ultimately tested by interventions on the replacement model, and interpretability
is ultimately judged by whether a reader can understand the final graph. Chapter 5
therefore evaluates these properties empirically. During construction, pruning
uses retained-mass thresholds on node and edge importance scores as a proxy for
importance retention. Clustering uses the atomicity loss $L_{\text{atom}}$ as a
proxy for coherent supernodes, and also uses the causal-preservation loss
$L_{\text{causal}}$ to discourage hiding strong causal edges inside supernodes.

\subsection{Pruning}
\label{sec:stage1}

Pruning reduces the attribution graph before clustering. This stage targets D1:
it keeps nodes and edges that lie on high-scoring paths between the seed
embedding nodes and seed logit nodes, and removes the rest so that later stages
operate on a tractable subgraph. The pruning rule first scores nodes by influence
and relevance, then scores edges on the node-pruned graph, and finally applies
retained-mass thresholds to both score sets.

\paragraph{Node influence and node relevance.}
Node influence measures how strongly a node can reach the seeded logits, while
node relevance measures how strongly a node is reached from the seeded input
tokens. Let
\[
W^{-}_{ij} = \frac{|W_{ij}|}{\sum_k |W_{ik}|},
\qquad
W^{+}_{ji} = \frac{|W_{ij}|}{\sum_k |W_{kj}|}
\]
be the backward and forward normalized absolute adjacencies, respectively.
Because $W_{ij}$ is the edge from source $j$ to target $i$, $W^{-}$ distributes a
target node's mass over its sources, while $W^{+}$ distributes a source node's
mass over its targets. The per-node influence and relevance vectors are
\begin{align}
\mathrm{Inf} &= \ell^\top \bigl(I - W^{-}\bigr)^{-1}
= \ell^\top \sum_{k \ge 0} (W^{-})^k
\;\in\; \mathbb{R}^N, \label{eq:inf}\\
\mathrm{Rel} &= \tau^\top \bigl(I - W^{+}\bigr)^{-1}
= \tau^\top \sum_{k \ge 0} (W^{+})^k
\;\in\; \mathbb{R}^N. \label{eq:rel}
\end{align}
Because $G$ is a DAG, both normalized adjacency matrices are nilpotent after a
topological ordering of the nodes. The Neumann series therefore terminates after
finitely many terms, so the inverse expressions are well-defined.

\paragraph{Node importance.}
A node is important only when both scores are high: it must be connected to the
seed input context and also influence the seed output. After normalizing
influence and relevance to $[0,1]$ across nodes, the node importance score is
\begin{equation}
\mathrm{Imp}_v =
\widehat{\mathrm{Inf}}_v^{\,\alpha}\,
\widehat{\mathrm{Rel}}_v^{\,1-\alpha},
\qquad \alpha \in [0,1],
\label{eq:importance}
\end{equation}
where $\widehat{\,\cdot\,}$ denotes the normalized score. The geometric mean is
used because it assigns zero importance whenever either factor is zero. This
property removes reconstruction-error nodes without a separate rule: an error node
is a source with no path from the seed embedding nodes, so its relevance is zero
and its combined importance is zero.

\paragraph{Node threshold.}
The node threshold $\eta_V \in [0,1]$ is a retained-mass threshold, not a raw
score cutoff. The implementation sorts node importance scores in descending
order, computes their cumulative fraction of total score mass, and finds the
score value at which the cumulative fraction first reaches $\eta_V$. It keeps
every node whose importance is at least that value. Embedding and logit nodes are
then retained unconditionally when endpoint preservation is enabled, because they
define the seed boundary of the summarization problem.

\paragraph{Edge influence and edge relevance.}
Edges are scored after node pruning, using the node-pruned adjacency
$\widetilde{W}$. Let $\widetilde{W}^{-}$ and $\widetilde{W}^{+}$ be the backward
and forward normalized absolute adjacencies computed from $\widetilde{W}$, and
let $\widetilde{\mathrm{Inf}}$ and $\widetilde{\mathrm{Rel}}$ be the corresponding
node influence and relevance scores on the node-pruned graph. For an edge from
source $j$ to target $i$, the edge influence and edge relevance are
\begin{align}
\mathrm{EInf}_{ij} &= \widetilde{W}^{-}_{ij}\,
\widetilde{\mathrm{Inf}}_i, \label{eq:einf}\\
\mathrm{ERel}_{ij} &= \widetilde{W}^{+}_{ji}\,
\widetilde{\mathrm{Rel}}_j. \label{eq:erel}
\end{align}
Edge influence assigns credit to an incoming edge according to the influence of
its target and the edge's normalized share of that target's incoming mass. Edge
relevance assigns credit according to the relevance of its source and the edge's
normalized share of that source's outgoing mass. The two quantities are flattened
over all matrix entries, normalized in one shared edge index space, and combined
with the same weighted geometric mean used for node importance.

\paragraph{Edge threshold and dangling-node removal.}
The edge threshold $\eta_E \in [0,1]$ has the same retained-mass meaning as the
node threshold. The implementation sorts the combined edge scores, finds the
score value where cumulative edge-score mass first reaches $\eta_E$, and keeps
all edges whose score is at least that value. Because edge pruning can disconnect
nodes that survived node pruning, the method then removes dangling nodes in a
fixed-point sweep. Embedding and error nodes must keep at least one outgoing
edge, logit nodes must keep at least one incoming edge, and feature nodes must
keep both an incoming and an outgoing edge. The sweep repeats until the node mask
is stable, yielding the pruned subgraph $(V',E')$.

\subsection{Clustering}
\label{sec:stage2}

Clustering turns the retained attribution subgraph into supernodes. It targets
D2, D3, and C2 at the same time: grouped features should have one causal role,
strong causal edges should remain visible between groups, and the number of groups
should fit the complexity budget. The stage fixes embedding and logit nodes as
singleton supernodes and clusters only the retained feature nodes
$V'_{\text{feature}} := V' \cap V_{\text{feature}}$.

\paragraph{Partition map.}
The clustering output is a partition map $f : V' \to V_H$. For every retained
embedding or logit node $v$, the map assigns $v$ to its singleton supernode
$\{v\}$. For feature nodes, the map may assign multiple nodes to the same
supernode. The number of resulting supernodes must satisfy $|V_H| \le K_{\max}$.

\paragraph{Role vectors.}
A feature's causal role is represented by its signed outgoing and incoming edge
profiles. For a retained feature node $u$, define
\begin{align}
(v_u^{\text{out}})_k &= W_{ku}\,\sqrt{\mathrm{Inf}_k},
\qquad k \in V', \nonumber\\
(v_u^{\text{in}})_k &= W_{uk}\,\sqrt{\mathrm{Rel}_k},
\qquad k \in V', \nonumber\\
r(u) &= \bigl[\,v_u^{\text{out}};\,v_u^{\text{in}}\,\bigr]
\;\in\; \mathbb{R}^{2|V'|}. \label{eq:role}
\end{align}
The outgoing half records which downstream nodes $u$ influences, weighted by the
influence of those targets. The incoming half records which upstream nodes
influence $u$, weighted by the relevance of those sources. The coordinates keep
the sign of $W$, so features that affect the same targets with opposite signs have
opposing role vectors rather than appearing similar.

\paragraph{Signed role similarity.}
The method compares two feature nodes by the signed cosine similarity of their
role vectors:
\begin{equation}
s_{ij} =
\cos\bigl(r(i), r(j)\bigr)
\;\in\; [-1,1],
\qquad i,j \in V'_{\text{feature}}.
\label{eq:sim}
\end{equation}
This similarity is positive when two features have aligned signed upstream and
downstream profiles, near zero when their profiles are unrelated, and negative
when their roles are antagonistic. The signed value is important for atomicity:
two features that drive the same downstream nodes with opposite signs should not
be merged merely because their edge supports overlap. Keeping the sign lets the
clustering objective reward aligned roles and penalize antagonistic roles in the
same expression.

\paragraph{Atomicity loss.}
Atomicity is measured by a signed correlation-clustering loss over the free
feature nodes. Let
$x_{ij} = \mathbf{1}[f(i) = f(j)]$ indicate whether feature nodes $i$ and $j$
share a supernode, and let $\theta \in [-1,1]$ be a resolution threshold. The
atomicity loss is
\begin{equation}
L_{\text{atom}}(f)
=
\sum_{\{i,j\} \subseteq V'_{\text{feature}}} x_{ij}\,(\theta - s_{ij}).
\label{eq:Latom}
\end{equation}
If $s_{ij} > \theta$, merging the pair contributes a negative term and is
rewarded. If $s_{ij} < \theta$, merging the pair contributes a positive term and
is penalized. Minimizing this loss therefore favours supernodes whose members
have aligned signed roles. In the experiments, $\theta = 0$ uses the sign of the
cosine as the merge boundary.

\paragraph{Causal-preservation loss.}
Atomicity alone can merge directly connected features, which hides their edge
inside a supernode. The causal-preservation loss measures the fraction of retained
edge mass absorbed internally by the partition:
\begin{equation}
L_{\text{causal}}(f)
=
\frac{\sum_{(u,v) \in E':\, f(u) = f(v)} |W_{vu}|}
     {\sum_{(u,v) \in E'} |W_{vu}|}
\;\in\; [0,1].
\label{eq:Lcausal}
\end{equation}
This term targets D3. It discourages merging pairs of nodes connected by large
retained edges, because such a merge would remove that causal relation from the
visible edge set of the summary graph.

\paragraph{Clustering objective and ILP.}
The partition is chosen by minimizing a weighted objective:
\begin{equation}
L(f) = L_{\text{atom}}(f) + \lambda\,L_{\text{causal}}(f),
\qquad \lambda \ge 0.
\label{eq:Lstage2}
\end{equation}
The parameter $\lambda$ controls the trade-off between grouping features with
similar roles and preserving explicit edges between them. A larger value of
$\lambda$ discourages merging directly connected features and tends to expose more
causal wiring in the final summary graph.

This objective is optimized as a binary integer linear program. The program uses
a binary variable $x_{ij}$ for each pair of retained feature nodes, where
$x_{ij}=1$ means that $i$ and $j$ are assigned to the same supernode. It also uses
binary representative variables $r_i$ to count the number of feature supernodes.
The objective is linear in the $x_{ij}$ variables: $L_{\text{atom}}$ is linear by
Eq.~\eqref{eq:Latom}, and the numerator of $L_{\text{causal}}$ is
$\sum_{(u,v)\in E'} |W_{vu}|\,x_{uv}$ for retained feature-feature edges.

The constraints ensure that the pair variables define a valid partition. Triangle
transitivity constraints such as
$x_{ij} + x_{jk} - x_{ik} \le 1$ prevent inconsistent assignments. The
representative variables count the number of clusters and enforce the budget
$K \le K_{\max}$. A layer-span constraint may additionally force
$x_{ij}=0$ for feature pairs whose layers are too far apart, preventing one
supernode from covering an overly broad section of the computation. For the
graphs considered in this thesis, the resulting mixed-integer program is solved
directly, so the clustering stage returns the optimizer of Eq.~\eqref{eq:Lstage2}
under the stated constraints.

\subsection{Interpreting}
\label{sec:stage3}

Interpreting converts the clustered pruned graph into the final summary graph.
This stage takes the partition $f$ from clustering, aggregates retained edges
between supernodes, removes cycles introduced by aggregation, and assigns
human-readable labels. It is responsible for enforcing the acyclicity constraint
C1.

\paragraph{Block-sum adjacency.}
For each ordered pair of distinct supernodes $(C,C')$, the signed block-sum edge
weight is
\begin{equation}
B_{C \to C'} =
\sum_{u \in C,\, v \in C'} W_{vu}\,\mathbf{1}[(u,v) \in E'].
\label{eq:block-sum}
\end{equation}
This quantity sums all retained attribution-graph edges from members of $C$ to
members of $C'$. The row-column convention is preserved: $W_{vu}$ is the edge
from source $u$ to target $v$. Diagonal blocks are omitted because edges inside a
supernode are represented as part of that node rather than as visible edges in
$H$. The off-diagonal block sums define a raw supernode adjacency, which may still
contain antiparallel edges or longer cycles.

\paragraph{Antiparallel collapse.}
The first acyclicity step removes two-node cycles. For each unordered pair
$\{C,C'\}$, let $a = B_{C \to C'}$ and $b = B_{C' \to C}$. The method keeps only
the net flow in the direction with larger magnitude:
\begin{equation}
(B_{C \to C'},\, B_{C' \to C}) \leftarrow
\begin{cases}
\bigl((|a|-|b|)\,\mathrm{sgn}(a),\,0\bigr) & \text{if } |a| > |b|, \\[2pt]
\bigl(0,\,(|b|-|a|)\,\mathrm{sgn}(b)\bigr) & \text{if } |b| > |a|, \\[2pt]
(0,\,0) & \text{if } |a| = |b|.
\end{cases}
\label{eq:antiparallel}
\end{equation}
This rule removes every antiparallel pair and leaves an edge only when one
direction has nonzero net dominance.

\paragraph{Back-edge removal.}
The second acyclicity step removes longer cycles by imposing a depth order on the
supernodes. Embedding supernodes receive depth $-\infty$, logit supernodes receive
depth $+\infty$, and each feature supernode receives the median layer of its
member features. Ties are broken by mean member layer, then by the smallest member
$\text{ctx\_idx}$, and finally by supernode index. After sorting by this order, the
method deletes every edge $C \to C'$ whose source ranks after its target. The
remaining adjacency is the summary adjacency $W_H$, and the final edge set is
\[
E_H = \{(C,C') : W_{H,\,C \to C'} \ne 0\}.
\]
By construction, all surviving edges point forward in the imposed order, so $H$ is
a DAG.

\paragraph{Labeling.}
Labeling is a display step rather than part of the graph construction objective.
Each supernode receives a short name and a coarse role, such as input, abstract,
or output, based on the feature descriptions of its members. These labels help the
reader inspect the final graph, but they do not affect $V_H$, $E_H$, $W_H$, or any
loss defined above.

\section{Chapter Summary}

This chapter formalized attribution-graph summarization and defined the proposed
three-stage method. The input is a signed directed acyclic attribution graph, and
the output is a directed acyclic summary graph whose supernodes partition a
retained subset of non-error nodes. The method first prunes the attribution graph
using node and edge influence--relevance scores with retained-mass thresholds. It
then clusters retained feature nodes with a signed role similarity and an ILP objective
$L_{\text{atom}} + \lambda L_{\text{causal}}$, which balances atomicity against
visible causal-edge preservation under a complexity budget. Finally, it aggregates
the clustered graph with block-sum edge weights, removes cycles, and labels the
resulting supernodes for inspection. The next chapter analyzes properties of this
construction, and Chapter 5 evaluates whether the resulting summary graphs are
faithful to the replacement model under steering interventions.

\end{document}
